\appendix

\section{Quality Control}
\label{app:qc}

\subsection{Human in the Task Synthesis Loops}
\label{app:qc-loop}

Each candidate task underwent a structured expert review before inclusion in
\datasetname. A domain-expert \textbf{Challenger} examined every task against
three criteria: whether the task description unambiguously covered all
verification checkpoints, whether every referenced entity and operation was
feasible in the deployed application environment, and whether the task
reflected a genuine professional workflow rather than a contrived exercise.
Issues were categorised by severity; any task with a critical or
high-severity finding was either returned for revision or rejected outright.
A senior \textbf{Refiner} then synthesised the critique, made the final
accept/revise/reject decision, and---for tasks requiring revision---specified
exactly what needed to change before the next review round. This
challenge--refine cycle repeated until the task met all criteria or was
discarded; many tasks underwent two or three rounds before acceptance.
Across all domains, only 45\% of candidate tasks survived the full review
process, reflecting the stringency applied to ensure that every included
task is both executable and professionally meaningful.

\subsection{Rubrics}
\label{app:qc-rubrics}

\paragraph{Text-only static check.}
For text-only tasks, we ask human expert reviewers to perform a static quality check according to the following rubric. Reviewers score each task along six dimensions on a 1--3 scale and mark three binary anti-pattern flags. This check is used to identify tasks that lack professional depth, contain unnatural cross-application transitions, have weak dependency structure, or are difficult to verify reliably.

\begin{tcolorbox}[
    title=Text-only Static Check Rubric,
    colback=gray!3,
    colframe=black!45,
    fonttitle=\bfseries,
    breakable
]
\textbf{Scoring scale.}
\begin{itemize}[]
    \item \textbf{1 = Poor}: clearly fails to satisfy the requirement.
    \item \textbf{2 = Fair}: partially satisfies the requirement but has clear weaknesses.
    \item \textbf{3 = Good}: fully satisfies the requirement.
\end{itemize}

\textbf{Six scoring dimensions.}

\begin{enumerate}
    \item \textbf{Professionalism.}
    \begin{itemize}
        \item \textbf{1}: Only requires generic UI operations and does not require domain knowledge.
        \item \textbf{2}: Involves some domain concepts, but only at a shallow level.
        \item \textbf{3}: Requires substantial domain knowledge, such as accounting rules, clinical workflows, HR policies, or event operations.
    \end{itemize}

    \item \textbf{Cross-app naturalness.}
    \begin{itemize}
        \item \textbf{1}: The involved applications have no real business relationship and are combined only to increase the number of apps.
        \item \textbf{2}: The relationship between applications is somewhat plausible, but some transitions lack clear motivation.
        \item \textbf{3}: Every application switch is naturally driven by business logic, and a real practitioner would plausibly follow the same workflow.
    \end{itemize}

    \item \textbf{Dependency depth.}
    \begin{itemize}
        \item \textbf{1}: The steps are independent and can be completed in any order.
        \item \textbf{2}: Some steps depend on earlier steps, but the dependency structure is weak.
        \item \textbf{3}: The task has a strict sequential dependency structure: outputs, data, computed values, or decisions from earlier steps determine the inputs of later steps.
    \end{itemize}

    \item \textbf{Verifiability.}
    \begin{itemize}
        \item \textbf{1}: The expected output is vague or subjective, such as ``write a reasonable document.''
        \item \textbf{2}: Some checkpoints have precise expected values, while others remain ambiguous.
        \item \textbf{3}: All checkpoints have explicit expected values and can be programmatically verified.
    \end{itemize}

    \item \textbf{Narrative coherence.}
    \begin{itemize}
        \item \textbf{1}: The task feels like a list of UI operations, with no unified business context.
        \item \textbf{2}: The task has a business background, but the narrative is not fully coherent.
        \item \textbf{3}: The task forms a single coherent business story, and every step contributes to the same overall objective.
    \end{itemize}

    \item \textbf{Complexity quality.}
    \begin{itemize}
        \item \textbf{1}: The high operation count mainly comes from repetitive data entry or over-specification, such as requiring exact wording.
        \item \textbf{2}: The task contains a mixture of genuine business complexity and redundant noise.
        \item \textbf{3}: The high operation count is driven by realistic business decisions and domain depth rather than artificial repetition.
    \end{itemize}
\end{enumerate}

\textbf{Three anti-pattern flags.}
Each flag is binary: \textbf{1} means the anti-pattern is present, and \textbf{0} means it is absent.

\begin{itemize}
    \item \textbf{{\ttfamily crm\_dumping\_ground}.}
    The final step merely creates a summary note in a CRM or document tool without meaningful business logic, and the tool is included only to add another application.

    \item \textbf{{\ttfamily parallel\_tasking}.}
    The steps across multiple applications have no data dependency. The operations in different applications can be completed independently or in parallel, and the order does not affect the result.

    \item \textbf{{\ttfamily spec\_overflow}.}
    The task description is overly detailed, such as specifying exact wording, cell values, or email body text. These specification details themselves become the main source of complexity rather than the underlying business logic.
\end{itemize}
\end{tcolorbox}

\paragraph{Multimodal static check.}
The multimodal static checker extends the text-only rubric to tasks that require images, PDFs, documents, audio, or video inputs. It keeps the six text-only dimensions but adds a seventh dimension, \textbf{multimodal feasibility}, so the total score becomes $7 \times 3 = 21$. This new dimension evaluates whether the referenced media actually exists, whether the media type can physically contain the requested information, and whether the task genuinely requires the agent to inspect the media rather than relying on values already exposed in the textual description. For tasks without multimodal input, this dimension is assigned a neutral score. The multimodal checker also introduces a new anti-pattern flag, {\ttfamily modal\_phantom}, which captures tasks that are only superficially multimodal. This flag is triggered when the referenced file path is missing or ambiguous, when the ground-truth value is leaked directly in the task description, or when the required field is impossible to infer from the referenced media type.

Compared with the text-only static check, the multimodal version also changes how cross-application workflows are judged. In text-only tasks, only application-to-application transitions count as cross-app transitions. In multimodal tasks, a media-to-application transition can also be a valid workflow boundary, such as extracting a value from a field photo, label image, invoice PDF, or document and entering it into a downstream SaaS application. To be considered high quality, such tasks must specify the media type, the exact field to extract, the downstream application, and the target field where the extracted information should be used.

\paragraph{Execution check.}
After static filtering, \datasetname further asks human expert reviewers to manually execute each task and assess its feasibility using the task-specific verifier and the resulting execution outcome. This stage is intended to assess task feasibility, verifier correctness, and failure attribution, rather than to evaluate model capability. For each task, reviewers inspect the task definition files, including {\ttfamily description.md}, {\ttfamily meta.json}, and {\ttfamily verify.py}, together with the human execution trace and verifier output. When necessary, reviewers also consult the involved applications' {\ttfamily APP\_SPEC.json}, user manuals, developer documentation, and container configuration to understand the intended UI workflow, terminology, data model, and API structure. The Execution Check Rubric below is the rubric used by human expert reviewers when manually executing and validating each task.

\begin{tcolorbox}[
    title=Execution Check Rubric,
    colback=gray!3,
    colframe=black!45,
    fonttitle=\bfseries,
    breakable
]
\textbf{Group A --- Alignment.}
This group checks the alignment between the task description and {\ttfamily verify.py}.
\begin{itemize}
    \item \textbf{A1.} Does the task goal in the description correspond to all checks in {\ttfamily verify.py}, with no missing or extra checks?
    \item \textbf{A2.} Are the concrete fields or values checked by {\ttfamily verify.py}, such as amounts, names, dates, or field names, clearly grounded in the task description?
    \item \textbf{A3.} Does the set of applications and operations involved in the task match the scope checked by {\ttfamily verify.py}?
\end{itemize}

\textbf{Group B --- Attribution.}
This group attributes the likely cause of any observed failure.
\begin{itemize}
    \item \textbf{B1.} If {\ttfamily verify.py} reports a failure, is the error possibly caused by ambiguity or mistakes in the task description rather than by expert operation? If all checks pass, this item is marked as false.
    \item \textbf{B2.} If {\ttfamily verify.py} reports a failure, is the error caused or constrained by limitations of the UI or the application? If all checks pass, this item is marked as false.
    \item \textbf{B3.} Is the overall human execution trace logically coherent, with clear progress between steps and no major jumps or contradictions?
\end{itemize}

\textbf{Group C --- Clarity.}
This group evaluates the clarity of the task description itself.
\begin{itemize}
    \item \textbf{C1.} Does the task description clearly specify the required operations or expected outcome, rather than giving a vague instruction?
    \item \textbf{C2.} Are the key parameters in the task description, such as amounts, names, dates, and field names, clear and unambiguous?
    \item \textbf{C3.} Does the task description avoid proprietary button names or UI labels that may not match the actual application interface?
\end{itemize}
\end{tcolorbox}

\section{Verification Methods}
\label{app:verify}

Tab.~\ref{tab:verify-types} provides a detailed breakdown of the three verification categories used in \datasetname and their subtypes.

\begin{table}[h]
\centering
\small
\caption{Verification methods used in \datasetname. Each checkpoint is assigned one of three categories depending on the nature of the expected output.}
\label{tab:verify-types}
\vspace{1em}
\renewcommand{\arraystretch}{1.25}
\begin{tabular}{llp{7.5cm}}
\toprule
\textbf{Category} & \textbf{Subtype} & \textbf{Description} \\
\midrule
\multirow{5}{*}{State-Check}
  & DB-Check & Queries the database to verify that expected records, fields, or relations exist. \\
  & API-Check & Calls backend APIs to confirm that the system state matches the expected outcome. \\
  & Numeric-Threshold-Check & Verifies that a numerical value falls within a specified tolerance range. \\
  & File-System-Check & Checks the existence, format, or content of files created during task execution. \\
  & Email/Message-Check & Verifies that expected emails or messages have been sent with correct recipients and content. \\
\midrule
\multirow{2}{*}{Content-Check}
  & Regex/String-Match & Matches extracted text against regular expressions or exact string patterns. \\
  & Doc-Extraction-Check & Extracts structured content from documents and verifies it against expected fields or keywords. \\
\midrule
\multirow{2}{*}{LLM-Judge}
  & LLM-Judge (text) & Uses an LLM to assess open-ended textual outputs such as reports, comments, or summaries. \\
  & LLM-Vision-Judge (image) & Uses a vision-language model to evaluate image-based outputs such as charts, screenshots, or visual artifacts. \\
\bottomrule
\end{tabular}
\end{table}

\section{Agent Action Space}
\label{app:actions}

Table~\ref{tab:actions} lists the 22 actions available to the \datasetname agent.

\begin{table}[h]
\centering
\small
\caption{The 22 actions available to the \datasetname agent. Browser UI actions interact with live web pages, file-system actions manage local artefacts (e.g.\ scratchpads, downloaded files), and \texttt{done} ends the episode. The \texttt{evaluate} (raw JavaScript) action is intentionally disabled to keep agent behaviour reproducible.}
\label{tab:actions}
\vspace{1em}
\setlength{\tabcolsep}{6pt}
\renewcommand{\arraystretch}{1.15}
\begin{tabular}{l l p{8cm}}
\toprule
\textbf{Category} & \textbf{Action} & \textbf{Description} \\
\midrule
\multirow{18}{*}{Browser UI}
  & \texttt{click}             & Click the indexed element on the current page. \\
  & \texttt{input}             & Type text into the indexed input/textarea (optionally clearing first). \\
  & \texttt{send\_keys}        & Send a raw keystroke or shortcut (e.g.\ \texttt{Enter}, \texttt{Ctrl+a}). \\
  & \texttt{scroll}            & Scroll the page (or an element) up or down by $N$ pages. \\
  & \texttt{select\_dropdown}  & Pick an option from a \texttt{<select>} dropdown by exact text. \\
  & \texttt{dropdown\_options} & Enumerate the options of an indexed dropdown. \\
  & \texttt{upload\_file}      & Upload a local file through an indexed file input. \\
  & \texttt{navigate}          & Open a URL (first-time entry into an application). \\
  & \texttt{go\_back}          & Navigate one step back in browser history. \\
  & \texttt{switch}            & Switch to another open tab by tab id. \\
  & \texttt{close}             & Close the tab identified by the given tab id. \\
  & \texttt{wait}              & Pause execution for $N$ seconds (e.g.\ to wait for loading). \\
  & \texttt{search}            & Issue a web search query (DuckDuckGo by default). \\
  & \texttt{extract}           & Use the LLM to extract structured information from the page. \\
  & \texttt{find\_elements}    & Locate elements by CSS selector and return their attributes. \\
  & \texttt{find\_text}        & Check whether a literal string appears on the page. \\
  & \texttt{search\_page}      & Regex/text search across visible page content. \\
  & \texttt{save\_as\_pdf}     & Save the current page as a PDF file. \\
\midrule
\multirow{3}{*}{File system}
  & \texttt{read\_file}        & Read a file from the agent's working directory. \\
  & \texttt{write\_file}       & Create or overwrite a file with the supplied content. \\
  & \texttt{replace\_file}     & Replace a substring inside an existing file. \\
\midrule
Completion
  & \texttt{done}              & Terminate the trajectory and report the final answer. \\
\bottomrule
\end{tabular}
\end{table}
